\section{实验与结果分析}

大语言模型评测需要结合任务场景、评价指标和输出可靠性设计评价口径\cite{luowen2024_llm_eval}。本章围绕基于大模型的课程知识实体识别与关系抽取系统，报告实验设计与运行结果。先介绍实验环境与配置参数，随后从课程语料构建、实体抽取评估、关系抽取评估、零样本与少样本抽取对比评估、知识图谱质量分析以及典型案例分析等角度，对系统运行结果进行评估与讨论。

\subsection{实验环境与配置}

本实验在统一的软硬件环境下运行，实验配置如表~\ref{tab:5-1}所示。硬件选用计算能力较强的服务器，软件环境基于Python 3.13，使用uv作为包管理器。系统核心依赖大语言模型API完成OCR识别、知识抽取与向量化表示三类任务。

\begin{table}[ht]
    \centering
    \caption{实验软硬件环境配置}
    \label{tab:5-1}
    \begin{tabular}{llp{8cm}}
    \toprule
    \textbf{配置类别} & \textbf{配置项} & \textbf{参数值} \\
    \midrule
    \multirow{3}{*}{硬件配置}
    & CPU & x86\_64架构多核处理器 \\
    & 内存 & 32GB DDR4 \\
    & 存储 & SSD固态硬盘 \\
    \midrule
    \multirow{3}{*}{软件环境}
    & 操作系统 & Linux (Ubuntu 24.04 LTS) \\
    & Python版本 & 3.13 \\
    & 包管理器 & uv \\
    \midrule
    \multirow{3}{*}{LLM配置}
    & OCR模型 & qwen3-vl-8b \\
    & 抽取模型 & deepseek-v4-flash \\
    & 嵌入模型 & qwen3-embedding-8b（向量维度1024） \\
    \bottomrule
    \end{tabular}
\end{table}

表~\ref{tab:5-2}列出了流水线核心参数配置。PDF渲染选用200 DPI分辨率以兼顾OCR识别质量；文本分块最大token数设为4000，平衡上下文完整性与处理效率；知识抽取采用few-shot模式提高准确性；同时启用严格证据验证机制，要求所有抽取得到的三元组均可溯源至原文。

\begin{table}[ht]
    \centering
    \caption{流水线核心参数配置}
    \label{tab:5-2}
    \begin{tabular}{llp{6cm}}
    \toprule
    \textbf{处理阶段} & \textbf{参数项} & \textbf{参数值} \\
    \midrule
    PDF转PNG & DPI分辨率 & 200 \\
    \midrule
    文本分块 & 最大token数 & 4000 \\
    & 分块策略 & 语义边界感知 \\
    \midrule
    知识抽取 & 抽取模式 & few-shot (fs) \\
    & 证据审查 & enabled \\
    & 严格证据 & enabled \\
    \midrule
    并发处理 & OCR并发数 & 10 \\
    & 抽取并发数 & 4 \\
    & 执行器类型 & ThreadPoolExecutor \\
    \bottomrule
    \end{tabular}
\end{table}

本实验使用ThreadPoolExecutor对OCR识别与知识抽取两个阶段进行并行加速。OCR阶段并发数设为10，知识抽取阶段并发数设为4，该配置在API速率限制与处理效率之间取得了平衡。

\subsection{课程语料构建结果}

本实验构建了覆盖计算机网络课程全内容的多元化语料库。根据\texttt{data/pdf}目录下的实际源文件统计，语料来源包括课程课件与讲义、RFC协议标准、实验与编程材料、作业与考核材料以及教材参考书，如表~\ref{tab:5-3}所示。

\begin{table}[H]
    \centering
    \caption{课程语料来源构成}
    \label{tab:5-3}
    \begin{tabular}{lp{8cm}r}
    \toprule
    \textbf{来源类型} & \textbf{具体内容} & \textbf{数量} \\
    \midrule
    课程课件与讲义 & 课程介绍、章节讲义、UnderHIT课件、Kurose配套章节课件等 & 115份 \\
    RFC/协议标准文档 & TCP/IP、HTTP、TLS、DNS、IPv6等相关RFC文档 & 72份 \\
    实验与编程材料 & Wireshark实验、Socket编程、实验指导书与FAQ等 & 44份 \\
    作业与考核材料 & Assignment、课后作业、历年试卷、模拟题、答案与解析等 & 33份 \\
    教材与参考书 & 谢希仁《计算机网络》、Kurose \& Ross《自顶向下方法》、《TCP/IP详解》系列及复习资料等 & 19份 \\
    \midrule
    \textbf{总计} & 205份PDF、69份PPT/PPTX、9份Markdown源文件 & \textbf{283份} \\
    \bottomrule
    \end{tabular}
\end{table}

语料库构建结果统计如表~\ref{tab:5-4}所示。在283份源文件中，205份PDF文档进入PDF渲染与OCR流程，并生成274个PNG输出目录（每一个目录对应一个PDF页码集合）。系统最终产生了16,438个语义文本块，总token数达10,398,922，平均每块632.6个token。

\begin{table}[ht]
    \centering
    \caption{语料库构建统计结果}
    \label{tab:5-4}
    \begin{tabular}{lr}
    \toprule
    \textbf{统计指标} & \textbf{数值} \\
    \midrule
    输入PDF文档数 & 205份 \\
    PNG输出目录数 & 274个 \\
    生成文本块数 & 16,438个 \\
    总token数 & 10,398,922 \\
    平均每块token数 & 632.6 \\
    \bottomrule
    \end{tabular}
\end{table}

按进入OCR流程的205份PDF文档计算，平均每份PDF产生约80个文本块。PPT/PPTX与Markdown源文件主要用于补充课程课件、实验说明和章节文本，教材与讲义的内容量较大。语料覆盖中英文双语材料，既有中文教材的系统阐述，也有英文RFC文档的原始技术规范，表明pipeline对多样化课程资源具备处理能力。

\subsection{实体抽取评估}

本实验从16,438个文本块中抽取了实体23,600个，实体类型分布如表~\ref{tab:5-5}所示。系统采用四分类实体模式：concept（概念）、mechanism（机制）、protocol（协议）、metric（性能指标）。

\begin{table}[ht]
    \centering
    \caption{实体类型分布统计}
    \label{tab:5-5}
    \begin{tabular}{lrrl}
    \toprule
    \textbf{实体类型} & \textbf{数量} & \textbf{占比} & \textbf{典型示例} \\
    \midrule
    概念(concept) & 16,719 & 70.8\% & 拥塞控制、流量控制、路由选择 \\
    机制(mechanism) & 3,063 & 13.0\% & 滑动窗口、CRC校验、三次握手 \\
    协议(protocol) & 2,189 & 9.3\% & TCP、UDP、HTTP、IP \\
    性能指标(metric) & 1,629 & 6.9\% & RTT、吞吐量、丢包率、带宽 \\
    \midrule
    \textbf{总计} & \textbf{23,600} & \textbf{100\%} & — \\
    \bottomrule
    \end{tabular}
\end{table}

表~\ref{tab:5-5}显示，概念类型实体占比最高(70.8\%)，这与课程材料以阐述抽象网络概念与原理为主的特性一致。机制类型实体占13.0\%，课程对实现技术给予了相当篇幅。协议类型实体仅占9.3\%，但协议在网络通信中处于基础地位，数量占比不反映其重要程度。性能指标类型实体占6.9\%，覆盖了网络性能评估的主要维度。

高频实体分布如表~\ref{tab:5-6}所示。TCP以2,657次的出现频率排在第一，UDP(1,411次)、IP(1,020次)位列其后，三者是网络协议栈的基础协议。

\begin{table}[ht]
    \centering
    \caption{高频实体统计（Top 10）}
    \label{tab:5-6}
    \begin{tabular}{lrclr}
    \toprule
    \textbf{排名} & \textbf{实体} & \textbf{频次} & \textbf{类型} & \textbf{占比*} \\
    \midrule
    1 & TCP & 2,657 & protocol & 11.3\% \\
    2 & UDP & 1,411 & protocol & 6.0\% \\
    3 & IP & 1,020 & protocol & 4.3\% \\
    4 & HTTP & 932 & protocol & 3.9\% \\
    5 & DNS & 723 & protocol & 3.1\% \\
    6 & RTT & 551 & metric & 2.3\% \\
    7 & IPv6 & 521 & protocol & 2.2\% \\
    8 & ICMP & 492 & protocol & 2.1\% \\
    9 & IPv4 & 469 & protocol & 2.0\% \\
    10 & OSPF & 451 & protocol & 1.9\% \\
    \bottomrule
    \multicolumn{5}{l}{\small *占总实体数的百分比}
    \end{tabular}
\end{table}

实体出现频次的范围为1至2,657，平均频次3.9次，呈典型的长尾分布。前10个高频实体累计占比约38\%，超过半数的实体(约12,000个)仅出现1～2次。该分布遵循幂律规律：高频术语在网络课程文献中被反复讨论，低频术语仅在特定上下文中出现。

TCP、UDP、IP等高频协议实体的频次远超均值，与它们在网络课程中的教学重心一致。系统识别了四类实体，为后续关系抽取与知识图谱构建提供了输入。

\subsection{关系抽取评估}

从23,600个实体中，本实验抽取了关系三元组44,526个。关系类型采用四分类体系：contains（包含）、applied\_to（应用于）、depends\_on（依赖于）、compared\_with（比较于），分布统计如表~\ref{tab:5-7}所示。

\begin{table}[ht]
    \centering
    \caption{关系类型分布统计}
    \label{tab:5-7}
    \begin{tabular}{lrr}
    \toprule
    \textbf{关系类型} & \textbf{数量} & \textbf{占比} \\
    \midrule
    包含(contains) & 15,251 & 34.3\% \\
    应用于(applied\_to) & 14,291 & 32.1\% \\
    依赖于(depends\_on) & 9,601 & 21.6\% \\
    比较于(compared\_with) & 5,383 & 12.1\% \\
    \midrule
    \textbf{总计} & \textbf{44,526} & \textbf{100\%} \\
    \bottomrule
    \end{tabular}
\end{table}

contains与applied\_to两类合计占66.4\%，反映了教育内容在知识结构上的两个特征：既强调概念的层级包含关系(如"TCP包含拥塞控制")，也关注机制的具体应用场景(如"滑动窗口应用于TCP")。depends\_on关系占21.6\%，源自网络协议栈的层次依赖特性(如"HTTP依赖于TCP")。

compared\_with关系占比最低(12.1\%)。比较分析属于高阶认知任务，在基础教材中出现频率较低，但这类关系表达了概念间的对比与区别，对完整描述课程知识结构是必要的。

高频三元组统计如表~\ref{tab:5-8}所示。

\begin{table}[ht]
    \centering
    \caption{高频关系三元组统计（Top 5）}
    \label{tab:5-8}
    \begin{tabular}{lcr}
    \toprule
    \textbf{关系三元组} & \textbf{频次} & \textbf{关系类型} \\
    \midrule
    TCP \textit{compared\_with} UDP & 421 & 比较 \\
    TCP \textit{contains} 拥塞控制 & 133 & 包含 \\
    IPv6 \textit{compared\_with} IPv4 & 130 & 比较 \\
    HTTP \textit{depends\_on} TCP & 122 & 依赖 \\
    TCP \textit{depends\_on} IP & 75 & 依赖 \\
    \bottomrule
    \end{tabular}
\end{table}

TCP与UDP之间的比较关系出现最为频繁：双向比较(TCP↔UDP)累计超过600次，说明传输层协议对比是课程反复涉及的话题。IPv6与IPv4的比较关系出现130次，与网络层协议演进的教学内容直接相关。

本实验启用了严格证据验证机制(strict\_evidence enabled)，所有三元组均附带可追溯的原文证据。该机制过滤了无依据的虚假关系，确保每条三元组均可回溯到源文本。

\subsection{零样本与少样本抽取对比评估}

为了让zero-shot（zs）与few-shot（fs）的抽取结果可以进行公平比较，本实验没有把任一模式的预测结果直接作为金标来源，而是以\texttt{data/chunks}中的原始文本块为回溯对象，构建并复核独立评估集。当前用于评估的金标文件为\texttt{data/eval/gold\_triples.json}，它与\texttt{data/eval/gold\_triples.500\_reviewed.json}内容一致，包含500条关系三元组和594个实体，覆盖88个源文件以及186个源文本块。核查过程中逐条检查了\texttt{source\_file}、\texttt{chunk\_index}与\texttt{evidence}字段：所有源文件均存在，所有chunk索引均有效，500条证据文本均为对应源文本块中的连续原文子串。其中有17条证据最初只是空白折叠后的文本，已替换为带原始换行的精确片段。

评估采用系统中的\texttt{kgeval}严格匹配口径。实体匹配要求实体名与类型完全一致，三元组匹配要求\texttt{head}、\texttt{head\_type}、\texttt{relation}、\texttt{tail}和\texttt{tail\_type}五个字段完全一致，不进行别名归一化、大小写折叠、后缀剥离或向量语义匹配。为了避免用全量图谱去对比500条抽样金标时把大量未标注但可能正确的关系误计为假阳性，实验先把四种raw抽取输出过滤到金标覆盖的186个源文本块，再在关闭embedding合并的条件下执行\texttt{kgmerge}，最后用同一份金标计算指标。该口径把比较范围限制在相同文本证据范围内，适合观察zs/fs提示方式和复核开关对严格命中的影响。

\begin{table}[ht]
    \centering
    \caption{零样本与少样本抽取在500条金标上的严格匹配结果}
    \label{tab:5-zsfs}
    \resizebox{\textwidth}{!}{
    \begin{tabular}{lrrrrrrrrr}
    \toprule
    \textbf{模式} & \textbf{Raw文件} & \textbf{Raw三元组} & \textbf{合并三元组} & \textbf{实体P} & \textbf{实体R} & \textbf{实体F1} & \textbf{三元组P} & \textbf{三元组R} & \textbf{三元组F1} \\
    \midrule
    \texttt{zs-nocheck} & 184 & 418 & 337 & 26.95\% & 51.68\% & 35.43\% & 18.99\% & 12.80\% & 15.29\% \\
    \texttt{zs-recheck} & 182 & 397 & 314 & 27.45\% & 51.35\% & 35.78\% & 17.83\% & 11.20\% & 13.76\% \\
    \texttt{fs-nocheck} & 171 & 325 & 290 & 24.36\% & 44.95\% & 31.60\% & 16.21\% & 9.40\% & 11.90\% \\
    \texttt{fs-recheck} & 176 & 392 & 304 & 26.04\% & 45.29\% & 33.07\% & 15.79\% & 9.60\% & 11.94\% \\
    \bottomrule
    \end{tabular}
    }
\end{table}

表~\ref{tab:5-zsfs}显示，在严格三元组匹配口径下，\texttt{zs-nocheck}的三元组F1最高（15.29\%），三元组精确率和召回率也高于另外三种设置；\texttt{zs-recheck}的实体F1最高（35.78\%）。在few-shot设置内部，启用复核后实体F1从31.60\%升至33.07\%，三元组召回率从9.40\%升至9.60\%，但三元组精确率从16.21\%降至15.79\%，说明复核在保留更多候选关系的同时，也引入了更多未被严格金标命中的三元组。

本实验中的金标是经人工复核的正例参考集，并非对186个源文本块中所有可能正确关系的穷尽标注。因此，表中的精确率应理解为保守的严格重合精确率，不等于完整语义精确率；召回率反映系统对500条已接受金标三元组的严格覆盖能力。在当前提示词和严格匹配条件下，zero-shot输出比few-shot更接近金标标注风格，后续优化应聚焦few-shot示例与金标标注规范的一致性，而非单纯增加示例数量。

为了直接衡量两种提示策略在抽取阶段的原始能力，本实验设计了\textbf{原始抽取层面}的逐块评估方案，使用专用的\texttt{kgevalraw}评测工具，在未经过合并的raw三元组上直接与金标进行逐块对比。评估维度涵盖实体层面F1（匹配键：name + type）、三元组层面F1（head、head\_type、relation、tail、tail\_type五字段严格匹配）和关系层面F1（仅匹配head\_type、relation、tail\_type），以及head\_type、tail\_type、relation三个属性的分类准确率。

表~\ref{tab:5-raw-zsfs}展示了zero-shot与few-shot在三个scope层面的精确率、召回率与F1值对比。

\begin{table}[ht]
    \centering
    \caption{原始抽取层面zero-shot与few-shot的scope级指标对比}
    \label{tab:5-raw-zsfs}
    \begin{tabular}{llrrrr}
    \toprule
    \textbf{Scope} & \textbf{Metric} & \textbf{ZS} & \textbf{FS} & \textbf{$\Delta$ (FS$-$ZS)} \\
    \midrule
    \multirow{3}{*}{Entity} & Precision & 40.00\% & 36.90\% & $-$3.10个百分点 \\
    & Recall    & 28.62\% & 27.27\% & $-$1.35个百分点 \\
    & \textbf{F1} & \textbf{33.37\%} & 31.36\% & $-$2.00个百分点 \\
    \midrule
    \multirow{3}{*}{Triple} & Precision & 16.44\% & 13.35\% & $-$3.09个百分点 \\
    & Recall    & 12.00\% & 9.80\% & $-$2.20个百分点 \\
    & \textbf{F1} & \textbf{13.87\%} & 11.30\% & $-$2.57个百分点 \\
    \midrule
    \multirow{3}{*}{Relation} & Precision & 55.00\% & 53.85\% & $-$1.15个百分点 \\
    & Recall    & 80.49\% & 68.29\% & $-$12.20个百分点 \\
    & \textbf{F1} & \textbf{65.35\%} & 60.22\% & $-$5.13个百分点 \\
    \bottomrule
    \end{tabular}
\end{table}

属性分类方面，ZS的head\_type准确率（86.67\%）与FS（84.93\%）接近，但FS在tail\_type上显著领先（89.04\% vs 75.24\%，+13.80个百分点），relation准确率两者相当（85.71\% vs 86.30\%）。

覆盖方面，ZS在186个金标块中找到182个对应raw文件（缺失4个），FS找到176个（缺失10个）。

实验结果显示：（1）Zero-shot在全部三个scope层面的F1指标上均优于few-shot（Entity F1: 33.37\% vs 31.36\%，Triple F1: 13.87\% vs 11.30\%，Relation F1: 65.35\% vs 60.22\%），与合并评估结论一致，说明zero-shot输出更接近金标标注风格。（2）Triple F1与Relation F1差距悬殊（零样本下65.35\% vs 13.87\%），源于实体名称严格匹配：当LLM使用原文完整表述（如"传输控制协议"）而金标使用缩写（"TCP"）时，大量语义正确但字符串不匹配的三元组被计为假阴性。（3）Few-shot仅在tail\_type分类上优于zero-shot（89.04\% vs 75.24\%，+13.80个百分点），但抽取缺失更多（10个vs 4个金标块找不到对应raw文件）。

\subsection{知识图谱质量分析}

在44,526个关系三元组的基础之上，本实验构建了包含23,600个节点、44,526条边的课程知识图谱。图谱的基本统计指标如表~\ref{tab:5-9}所示。

\begin{table}[ht]
    \centering
    \caption{知识图谱基本统计指标}
    \label{tab:5-9}
    \begin{tabular}{lr}
    \toprule
    \textbf{统计指标} & \textbf{数值} \\
    \midrule
    实体节点数 & 23,600 \\
    关系边数 & 44,526 \\
    平均节点度数 & 3.77 \\
    实体平均出现频次 & 3.9次 \\
    高频实体占比(Top 10) & 38\% \\
    \bottomrule
    \end{tabular}
\end{table}

平均节点度数约3.8，与课程知识中高频术语多连接、低频术语少连接的组织规律一致。系统通过词汇级别名表和语义级嵌入向量双层归一化处理实体名称变体，并以TCP等高频协议为枢纽连接多个子概念与关系。

\subsection{案例分析}

本节选取两个典型子图进行分析。TCP协议子图包含对比关系（TCP \textit{compared\_with} UDP，双向616次）、包含关系（TCP \textit{contains} 拥塞控制133次、流量控制89次）和依赖关系（HTTP \textit{depends\_on} TCP 122次、TCP \textit{depends\_on} IP 75次），覆盖协议对比、内部机制和协议栈依赖三个知识维度。协议演进对比子图展示了IPv6 \textit{compared\_with} IPv4（130次）等课程核心对比主题，系统通过统计频次有效识别出教学重点。


\subsection{本章小结}

本章对系统进行了实验评估。系统构建出包含23,600个实体和44,526个关系三元组的课程知识图谱，实体分布以概念类型为主（70.8\%），关系分布以contains和applied\_to为主（66.4\%），平均节点度数3.8。零样本与少样本对比表明，zero-shot在当前提示词与严格匹配条件下整体优于few-shot（合并层面三元组F1: 15.29\% vs 11.94\%，原始层面: 13.87\% vs 11.30\%），few-shot在tail\_type分类上具有局部优势（89.04\% vs 75.24\%）。
